\chapter{Covariance and Correlation}

A joint distribution contains everything there is to know about two random
variables, but often we want a single summary number: when one variable goes up,
does the other tend to go up as well?

This chapter develops the two standard measures of linear association:
\textbf{covariance} and the scaled \textbf{correlation coefficient}. We examine
why correlation lies in $[-1,1]$, why zero correlation does not imply
independence, and how covariance matrices describe higher-dimensional data.

\section{The Question Covariance Answers}
\label{ch14:question}

\begin{intuitionbox}
Variance asks: \emph{how much does a single quantity wiggle around its mean?}
Covariance asks the natural two-variable sequel: \emph{when one quantity wiggles
above its mean, does the other tend to wiggle above its mean too, or below it?}
Covariance is, almost literally, "co-variance" -- variance shared between two
things.
\end{intuitionbox}

Let us make the story concrete. Imagine we sample students at random and record two
numbers for each: $X$, the hours they slept the night before an exam, and $Y$,
their exam score. Each student is a dot in the plane. If well-rested students tend
to score higher, the dots will lean upward to the right. If sleep makes no
difference, the dots will form a shapeless blob. If somehow sleep \emph{hurt}
scores, the dots would lean downward.

The trick that turns this picture into a number is to stop looking at the raw values
$X$ and $Y$ and instead look at the \emph{deviations} from their means,
\[
X - \mu_X, \qquad Y - \mu_Y, \qquad \text{where } \mu_X = \E[X],\ \mu_Y = \E[Y].
\]
A deviation is positive when the value is above average and negative when it is
below. Now multiply the two deviations together for each student and ask what that
product tends to be.

\begin{itemize}
\item A well-rested, high-scoring student has $X - \mu_X > 0$ and $Y - \mu_Y > 0$,
so the product $(X-\mu_X)(Y-\mu_Y)$ is \emph{positive}.
\item A sleep-deprived, low-scoring student has $X - \mu_X < 0$ and
$Y - \mu_Y < 0$: a negative times a negative is again \emph{positive}.
\item A student who is rested but scores low (or tired but scores high) contributes
a \emph{negative} product.
\end{itemize}

So the sign of the product $(X-\mu_X)(Y-\mu_Y)$ records whether a given point sits
in the "agreeing" quadrants (both above average or both below) or the "disagreeing"
quadrants. The \emph{average} of that product over the whole population is the
covariance.

\begin{definition}[Covariance]
\label{ch14:def-cov}
For random variables $X$ and $Y$ with finite variances, the \textbf{covariance} of
$X$ and $Y$ is
\[
\Cov(X, Y) \;=\; \E\!\big[(X - \mu_X)(Y - \mu_Y)\big],
\qquad \mu_X = \E[X],\ \mu_Y = \E[Y].
\]
\end{definition}

\begin{figure}[H]
\centering
\begin{tikzpicture}
\begin{axis}[
  width=10cm, height=7cm,
  xlabel={$X - \mu_X$},
  xlabel style={at={(axis cs:3.1,-0.2)}, anchor=north east, font=\small},
  xmin=-3.2, xmax=3.2, ymin=-3.2, ymax=3.2,
  axis lines=middle, xtick=\empty, ytick=\empty,
  enlargelimits=false, clip=false,
]
% axis label for Y
\node[anchor=south east, font=\small] at (axis cs:-0.15,3.1) {$Y - \mu_Y$};
% quadrant shading labels
\node[targetgreen!80!black, anchor=west] at (axis cs:0.4,2.5) {\small $(+)(+) = +$};
\node[targetgreen!80!black, anchor=east] at (axis cs:-0.3,-2.6) {\small $(-)(-) = +$};
\node[trapred] at (axis cs:-2.0,2.2) {\small $(-)(+) = -$};
\node[trapred] at (axis cs:2.0,-2.2) {\small $(+)(-) = -$};
% scatter leaning upward (positive covariance)
\addplot[only marks, mark=*, mark size=1.2pt, draw=sourceblue, fill=sourceblue]
  table {
  -2.4 -2.1
  -1.9 -1.4
  -1.6 -2.0
  -1.1 -0.6
  -0.8 -1.3
  -0.5 0.1
  -0.3 -0.7
  0.0 0.4
  0.2 -0.2
  0.5 0.9
  0.8 0.3
  1.1 1.6
  1.4 0.7
  1.7 1.9
  2.0 1.3
  2.3 2.4
  -1.3 0.4
  0.9 -0.5
  1.6 2.2
  -2.0 -1.0
  };
\end{axis}
\end{tikzpicture}
\caption{Deviations from the means split the plane into four quadrants. Points in
the upper-right and lower-left contribute positive products; points in the other
two contribute negative products. Covariance is the average product. This cloud
leans upward, so the positive quadrants win and the covariance is positive.}
\label{ch14:fig-quadrants}
\end{figure}

\begin{happeningbox}
Covariance measures the balance of deviation products between quadrants. Each point
contributes $(X - \mu_X)(Y - \mu_Y)$: positive if both coordinates deviate from their
means in the same direction, and negative if they deviate in opposite directions,
scaled by the distance from the mean axes. Points far from the means have a large
effect on the average, while points close to the means contribute little. A point
exactly on either mean line ($X = \mu_X$ or $Y = \mu_Y$) contributes zero,
because one of its deviations is zero. When a distribution is symmetric across
the mean axes, the positive and negative quadrant contributions cancel to zero.
\end{happeningbox}

\subsection{A computational shortcut}

The definition in terms of deviations is the one to \emph{think} with, but it is
rarely the easiest one to \emph{compute} with. Expanding the product gives a formula
that needs only the three expectations $\E[XY]$, $\E[X]$, $\E[Y]$.

\begin{proposition}[Shortcut formula]
\label{ch14:cov-shortcut}
\[
\Cov(X, Y) \;=\; \E[XY] - \E[X]\,\E[Y].
\]
\end{proposition}

\begin{whybox}
Expand the product inside the expectation and use linearity. Writing
$\mu_X = \E[X]$ and $\mu_Y = \E[Y]$ (these are constants),
\begin{align*}
\Cov(X,Y)
&= \E\big[(X-\mu_X)(Y-\mu_Y)\big] \\
&= \E\big[XY - \mu_X Y - \mu_Y X + \mu_X \mu_Y\big] \\
&= \E[XY] - \mu_X \E[Y] - \mu_Y \E[X] + \mu_X \mu_Y \\
&= \E[XY] - \mu_X \mu_Y - \mu_Y \mu_X + \mu_X \mu_Y \\
&= \E[XY] - \E[X]\,\E[Y].
\end{align*}
The two middle terms collapsed because $\E[\mu_X Y] = \mu_X \E[Y] = \mu_X \mu_Y$,
and similarly for the other. Notice the family resemblance to the variance
shortcut $\Var(X) = \E[X^2] - (\E[X])^2$ -- indeed, setting $Y = X$ recovers it
exactly, since $\Cov(X, X) = \E[X^2] - (\E[X])^2 = \Var(X)$.
\end{whybox}

\begin{focusbox}
Covariance is the \emph{generalization of variance to two variables}. The single
most useful identity to keep in your head is
\[
\Cov(X, X) = \Var(X).
\]
Variance is just a variable's covariance with itself. Everything we know about
variance is a special case of what we are about to learn about covariance.
\end{focusbox}

\begin{recoverybox}
\textbf{Pause and think:} If you flip a fair coin, let $X = 1$ for heads ($0$ for tails), and set $Y = 1 - X$ (the tails indicator). What is $\Cov(X, Y)$ without writing out any sums?

Because $X + Y = 1$ is constant, whenever $X$ rises above its mean, $Y$ drops below its mean by the exact same amount. The two deviations always have opposite signs, so the product is always negative (or zero). Using self-covariance:
\[
\Cov(X, Y) = \Cov(X, 1 - X) = \Cov(X, -X) = -\Cov(X, X) = -\Var(X) = -\tfrac{1}{4}.
\]
They move in perfect lockstep opposition.
\end{recoverybox}

\subsection{First worked example: two dice}

\begin{example}[Sum and first die]
\label{ch14:ex-dice}
Roll two fair, independent dice. Let $D_1$ be the first die's value and let
$S = D_1 + D_2$ be the sum. Intuitively, when the first die is large the sum tends
to be large, so we expect a positive covariance. Let us compute $\Cov(D_1, S)$.

Bilinearity (which just means that covariance is linear in each input separately, proved formally in Section~\ref{ch14:properties}) is the fast route,
but doing it by hand first builds trust. We have $\E[D_1] = 3.5$ and $\E[S] = 7$.
For the cross term,
\[
\E[D_1 S] = \E[D_1(D_1 + D_2)] = \E[D_1^2] + \E[D_1 D_2].
\]
Now $\E[D_1^2] = \tfrac{1}{6}(1 + 4 + 9 + 16 + 25 + 36) = \tfrac{91}{6}$, and
because the dice are independent, $\E[D_1 D_2] = \E[D_1]\E[D_2] = 3.5 \times 3.5
= 12.25$. So
\[
\E[D_1 S] = \tfrac{91}{6} + 12.25 = 15.1\overline{6} + 12.25 = 27.41\overline{6}.
\]
Therefore
\[
\Cov(D_1, S) = \E[D_1 S] - \E[D_1]\E[S]
= 27.41\overline{6} - 3.5 \times 7
= 27.41\overline{6} - 24.5 = 2.91\overline{6} = \tfrac{35}{12}.
\]
Positive, as our intuition demanded. And note $\tfrac{35}{12} = \Var(D_1)$: the
covariance of a die with the sum equals the variance of that one die, because the
\emph{other} die contributes nothing to the covariance (it is independent of
$D_1$). That is a preview of the bilinearity rules.
\end{example}

\section{The Algebra of Covariance}
\label{ch14:properties}

Covariance obeys a small set of rules that make it pleasant to manipulate. They all
follow from the linearity of expectation, and they will pay for themselves many
times over -- especially when we compute the variance of a sum.

\begin{proposition}[Properties of covariance]
\label{ch14:cov-props}
For random variables $X, Y, Z$ and constants $a, b, c$:
\begin{enumerate}
\item \textbf{Symmetry:} $\Cov(X, Y) = \Cov(Y, X)$.
\item \textbf{Self-covariance:} $\Cov(X, X) = \Var(X)$.
\item \textbf{Constants don't covary:} $\Cov(X, c) = 0$.
\item \textbf{Pulling out scalars and shifts:}
$\Cov(aX + b,\; cY + d) = ac\,\Cov(X, Y)$.
\item \textbf{Bilinearity:}
$\Cov(aX + bY,\; Z) = a\,\Cov(X, Z) + b\,\Cov(Y, Z)$,
and likewise in the second slot.
\end{enumerate}
\end{proposition}

\begin{whybox}
Symmetry is obvious from the definition (multiplication commutes). Self-covariance
we already saw. For (3), a constant has zero deviation: $c - \E[c] = 0$, so the
product is identically zero. For (4), adding constants $b, d$ shifts the means by
the same amounts, so the deviations $X - \mu_X$ and $Y - \mu_Y$ are unchanged by
$b$ and $d$, while the scalars $a, c$ factor straight out of the product. For
bilinearity, deviations add: the deviation of $aX + bY$ is
$a(X-\mu_X) + b(Y-\mu_Y)$, and expectation is linear, so
\[
\Cov(aX+bY, Z) = \E\big[(a(X-\mu_X)+b(Y-\mu_Y))(Z-\mu_Z)\big]
= a\Cov(X,Z) + b\Cov(Y,Z). \qedhere
\]
\end{whybox}

\begin{intuitionbox}
Property (4) carries a central lesson. The shifts $b$ and $d$ \emph{vanish} -- moving
your data left/right or up/down cannot change how two variables co-move. But the
scales $a$ and $c$ \emph{survive and multiply in}. This means covariance is
sensitive to \emph{units}. Measure a person's height in centimeters instead of
meters and every covariance involving height multiplies by $100$. Covariance has a
unit problem, and that problem is exactly what correlation will fix.
\end{intuitionbox}

\subsection{The payoff: variance of a sum}

Bilinearity immediately simplifies the computation of the variance of a sum.
Back in the variance chapter we showed that variances add for independent
variables, but noted that the general case requires covariance. With bilinearity
in hand, the general formula follows directly.

\begin{theorem}[Variance of a sum]
\label{ch14:var-sum}
For any random variables $X$ and $Y$ with finite variance,
\[
\Var(X + Y) = \Var(X) + \Var(Y) + 2\,\Cov(X, Y).
\]
More generally, for $X_1, \dots, X_n$,
\[
\Var\!\Big(\sum_{i=1}^n X_i\Big)
= \sum_{i=1}^n \Var(X_i) + 2\!\!\sum_{1 \le i < j \le n}\!\! \Cov(X_i, X_j)
= \sum_{i=1}^n \sum_{j=1}^n \Cov(X_i, X_j).
\]
\end{theorem}

\begin{whybox}
Just expand using $\Var(W) = \Cov(W, W)$ and bilinearity:
\[
\Var(X+Y) = \Cov(X+Y, X+Y)
= \Cov(X,X) + \Cov(X,Y) + \Cov(Y,X) + \Cov(Y,Y),
\]
and the two cross terms are equal by symmetry, giving the factor of $2$. For the
$n$-term version, $\Cov\big(\sum_i X_i, \sum_j X_j\big)$ expands into the full
double sum of all pairwise covariances; the diagonal terms $i = j$ are the
variances, and the off-diagonal terms come in symmetric pairs.
\end{whybox}

\begin{aestheticbox}
This identity is central throughout probability and statistics.
It says: \emph{the variance of a total is the sum of the variances plus a
correction for how the parts move together}. If the parts are uncorrelated, the
correction vanishes and variances simply add -- which is why, in the chapters
ahead, the standard error of an average shrinks like $1/\sqrt{n}$ when samples are
independent. When the parts are positively correlated, totals are \emph{more}
volatile than you would guess from the pieces alone; this is the mathematics behind
why a poorly diversified portfolio of correlated stocks is so risky, and why
diversification (mixing weakly correlated assets) reduces risk. We will return to
this in the application box below.
\end{aestheticbox}

\begin{example}[Variance of a binomial, the slick way]
\label{ch14:ex-binom}
Let $X \sim \Bin(n, p)$. Write $X = \sum_{i=1}^n I_i$ where $I_i$ is the indicator
of a success on trial $i$. The trials are independent, so $\Cov(I_i, I_j) = 0$ for
$i \ne j$ (we will prove independence kills covariance in a moment). Hence all the
cross terms in Theorem~\ref{ch14:var-sum} vanish and
\[
\Var(X) = \sum_{i=1}^n \Var(I_i) = \sum_{i=1}^n p(1-p) = np(1-p).
\]
No messy sums over the binomial pmf -- the covariance bookkeeping did the work.
\end{example}

\section{Independence, Covariance, and a Famous Trap}
\label{ch14:indep}

\begin{theorem}[Independence implies zero covariance]
\label{ch14:indep-implies-zero}
If $X$ and $Y$ are independent (and have finite variances), then
$\Cov(X, Y) = 0$.
\end{theorem}

\begin{whybox}
Independence means the joint factors: $\E[XY] = \E[X]\,\E[Y]$ (this is the
factorization of expectation for independent variables, proved in the
joint-distributions chapter). Plug this straight into the shortcut formula:
\[
\Cov(X, Y) = \E[XY] - \E[X]\E[Y] = \E[X]\E[Y] - \E[X]\E[Y] = 0. \qedhere
\]
\end{whybox}

When $\Cov(X, Y) = 0$ we call $X$ and $Y$ \textbf{uncorrelated}. So independence is
a \emph{sufficient} condition for being uncorrelated. The burning question is
whether it is \emph{necessary} -- does uncorrelated imply independent? The answer is
a flat \textbf{no}, and the reason rewards a careful look, because the confusion
is one of the most common in applied probability.

\begin{trapbox}
\textbf{Uncorrelated does NOT mean independent.} Covariance only detects
\emph{linear} co-movement -- the straight-line tendency for one variable to rise as
the other rises. A relationship can be ironclad and deterministic yet have zero
covariance, simply because it is not a \emph{straight-line} relationship. Zero
covariance rules out a linear trend; it rules out nothing else.
\end{trapbox}

\begin{example}[The U-shape that fools covariance]
\label{ch14:ex-ushape}
Let $X$ be uniform on $\{-2, -1, 0, 1, 2\}$ and let $Y = X^2$. Then $Y$ is a
\emph{perfect, deterministic function} of $X$ -- knowing $X$ tells you $Y$ exactly,
so they are about as far from independent as two variables can be. Yet
\[
\E[X] = 0, \qquad
\E[XY] = \E[X^3] = \tfrac{1}{5}\big((-2)^3 + (-1)^3 + 0 + 1^3 + 2^3\big) = 0,
\]
so $\Cov(X, Y) = \E[XY] - \E[X]\E[Y] = 0 - 0 = 0$. The covariance is exactly zero.
What happened? The relationship is a symmetric U (a parabola): as $X$ goes from
negative to positive, $Y$ first falls then rises. The rising right half exactly
cancels the falling left half, so covariance -- which measures only overall
linear tilt -- registers zero.
\end{example}

\begin{figure}[H]
\centering
\begin{tikzpicture}
\begin{axis}[
  width=10cm, height=6cm,
  xlabel={$X$}, ylabel={$Y = X^2$},
  xmin=-2.6, xmax=2.6, ymin=-0.5, ymax=4.8,
  axis lines=middle, samples=60, domain=-2.3:2.3,
  xtick={-2,-1,0,1,2}, ytick={0,1,2,3,4},
  enlargelimits=false,
]
\addplot[densityred, very thick, domain=-2.3:2.3] {x^2};
\addplot[only marks, mark=*, mark size=2.4pt, draw=sourceblue, fill=sourceblue]
  coordinates {(-2,4)(-1,1)(0,0)(1,1)(2,4)};
\draw[gray, dashed] (axis cs:0,-0.5) -- (axis cs:0,4.8);
\end{axis}
\end{tikzpicture}
\caption{$Y = X^2$ with $X$ symmetric about $0$. The dependence is total, but it is
purely curved -- there is no linear tilt -- so $\Cov(X, Y) = 0$. Covariance is
blind to symmetric, nonlinear relationships.}
\label{ch14:fig-ushape}
\end{figure}

\begin{happeningbox}
The deep reason for the trap: covariance is a \emph{single number}, and a single
number can only summarize the \emph{first-order, straight-line} part of a
relationship. Think of it as the best you could do if you were forced to predict
$Y$ from $X$ using a straight line. If the true relationship is a symmetric curve,
the best straight line is flat, and a flat line means zero covariance -- even though
$X$ and $Y$ are locked together. The lesson generalizes: \emph{the absence of a
linear signal is not the absence of a relationship.} Always plot your data; a
scatterplot sees the U-shape that the covariance number erases.
\end{happeningbox}

\begin{recoverybox}
\textbf{Pause and think:} Let $X \sim \Unif(-1, 1)$ and set $Y = |X|$. Are $X$ and $Y$ independent? What is $\Cov(X, Y)$?

Knowing $X$ completely determines $Y$, so they are strictly dependent. Yet $\E[X] = 0$, and by symmetry $\E[XY] = \E[X|X|] = \int_{-1}^1 x|x|\,\tfrac{1}{2}\,\dd x = 0$ because the integrand is an odd function (it flips sign across zero). Hence $\Cov(X, Y) = 0 - 0 = 0$. Just like the parabola $Y = X^2$, the V-shaped curve folds symmetric positive and negative deviations directly onto each other, creating zero linear covariance despite an absolute deterministic relationship.
\end{recoverybox}

There is exactly one important setting where uncorrelated \emph{does} upgrade to
independent: when the pair $(X, Y)$ is jointly Gaussian (bivariate normal). For
jointly normal variables, zero correlation is equivalent to independence. This is a
special feature of the normal distribution, not a general truth, and we will meet it
again when we study the multivariate normal.

\section{Standardizing: The Correlation Coefficient}
\label{ch14:correlation}

Covariance has the unit problem we flagged earlier. Its numerical value depends on
the scales of $X$ and $Y$, so a covariance of $2.9$ (as in our two-dice example)
is meaningless in isolation -- is that a strong relationship or a weak one? We
cannot tell without knowing how spread out the variables are. The fix is to
\emph{standardize}: divide out each variable's own spread, measured by its standard
deviation. The result is the \textbf{correlation coefficient}.

\begin{definition}[Correlation coefficient]
\label{ch14:def-corr}
For random variables $X$ and $Y$ with $\SD(X), \SD(Y) > 0$, the \textbf{Pearson
correlation coefficient} is
\[
\Corr(X, Y) \;=\; \rho_{X,Y} \;=\;
\frac{\Cov(X, Y)}{\SD(X)\,\SD(Y)}
\;=\;
\frac{\Cov(X, Y)}{\sqrt{\Var(X)\,\Var(Y)}}.
\]
\end{definition}

\begin{intuitionbox}
The cleanest way to see what correlation is: first \emph{standardize} each variable
into a z-score by subtracting its mean and dividing by its standard deviation,
\[
\tilde X = \frac{X - \mu_X}{\SD(X)}, \qquad
\tilde Y = \frac{Y - \mu_Y}{\SD(Y)}.
\]
Each $\tilde X$, $\tilde Y$ now has mean $0$ and variance $1$ -- pure, unit-free
numbers stripped of their original scale. Then
\[
\Corr(X, Y) = \Cov(\tilde X, \tilde Y) = \E[\tilde X \tilde Y].
\]
Correlation is just the covariance of the \emph{standardized} variables. It asks
whether deviations align, just as covariance does, but with both variables normalized
by their standard deviations so that the result is dimensionless and directly comparable
across different pairs of variables.
\end{intuitionbox}

Because $\Corr$ divides covariance by exactly the scale factors that covariance was
sensitive to, correlation is \textbf{scale- and shift-invariant}: for positive
constants $a, c$ and any $b, d$,
\[
\Corr(aX + b,\; cY + d) = \Corr(X, Y).
\]
(If exactly one of $a, c$ is negative, the correlation flips sign; if both are negative, the signs cancel and $\Corr$ is unchanged, because each negative scale
reverses the direction of co-movement.) Switching from meters to centimeters, or
from Celsius to Fahrenheit, leaves correlation untouched. That invariance is
precisely why correlation -- not covariance -- is the number people actually report.

\begin{recoverybox}
\textbf{Pause and think:} If $\Corr(X, Y) = 0.8$, what is $\Corr(3X + 5, -2Y + 4)$?

The additive shifts $+5$ and $+4$ do not affect deviations from the mean. Scaling $X$ by $+3$ changes its unit spread without changing its orientation. But multiplying $Y$ by $-2$ flips every deviation across its mean: what was above average becomes below average. The positive co-movement turns into negative co-movement, giving exactly $-0.8$.
\end{recoverybox}

\section{Why Correlation Lives in $[-1, 1]$}
\label{ch14:bounds}

The central fact about correlation is that it is bounded:
\[
-1 \;\le\; \Corr(X, Y) \;\le\; 1.
\]
A correlation of $+1$ means perfect positive linear agreement; $-1$ means perfect
negative linear agreement; $0$ means no linear relationship. Values in between
interpolate. But \emph{why} can it never escape this interval? Why not a correlation
of $1.4$? There are two ways to see it.

\subsection{The elementary proof: a variance can't be negative}

\begin{theorem}[Correlation bound]
\label{ch14:corr-bound}
For any $X, Y$ with positive, finite standard deviations,
$\,-1 \le \Corr(X, Y) \le 1$. The boundary values $\Corr(X, Y) = \pm 1$ occur if
and only if $Y$ is an exact linear function of $X$, that is, $Y = aX + b$ for
constants $a \ne 0$ and $b$ (with $a > 0$ for $+1$ and $a < 0$ for $-1$).
\end{theorem}

\begin{whybox}
Standardize as above, so $\tilde X, \tilde Y$ each have variance $1$, and let
$\rho = \Corr(X,Y) = \E[\tilde X \tilde Y]$. Now consider the variance of the
\emph{sum} and the \emph{difference} of the standardized variables -- and remember
that a variance can never be negative. Using
Theorem~\ref{ch14:var-sum} and $\Var(\tilde X) = \Var(\tilde Y) = 1$,
\[
0 \le \Var(\tilde X + \tilde Y) = 1 + 1 + 2\rho = 2(1 + \rho)
\;\Longrightarrow\; \rho \ge -1,
\]
\[
0 \le \Var(\tilde X - \tilde Y) = 1 + 1 - 2\rho = 2(1 - \rho)
\;\Longrightarrow\; \rho \le 1.
\]
Together these pin $\rho$ inside $[-1, 1]$. And equality is forced exactly when one
of those variances is \emph{zero}. A variable with zero variance is constant, so
$\rho = 1$ means $\tilde X - \tilde Y$ is constant -- i.e. $\tilde Y = \tilde X +
\text{const}$, which unwinds to $Y = aX + b$ with $a > 0$. Likewise $\rho = -1$
means $\tilde X + \tilde Y$ is constant, giving a perfect line with negative slope.
\end{whybox}

\begin{happeningbox}
That proof repays a second look, because it explains \emph{what the boundary means}.
Correlation hits $\pm 1$ precisely when the points stop being a cloud and collapse
onto a single straight line -- because that is the only way a variance can shrink to
zero. The closer $\abs{\rho}$ is to $1$, the tighter the cloud hugs some line. So
$\rho$ is best read as a measure of \emph{how close the scatter is to lying on a
straight line}, with the sign telling you which way the line tilts.
\end{happeningbox}

\subsection{The structural proof: Cauchy--Schwarz}

The bound is the probabilistic face of one of the
most fundamental inequalities in mathematics: the Cauchy--Schwarz inequality. Seeing
this connection turns the bound from a computational curiosity into a piece of
geometry.

\begin{connectionbox}
\textbf{Covariance is an inner product; correlation is the cosine of an angle.}
An inner product is the algebraic machinery for measuring lengths and angles. In ordinary Euclidean geometry, the dot product $\mathbf{u} \cdot \mathbf{v}$ tells you how much two arrows align. Recall from linear algebra that for vectors $\mathbf{u}, \mathbf{v}$, the dot product satisfies the \emph{Cauchy--Schwarz inequality}
\[
\abs{\langle \mathbf{u}, \mathbf{v}\rangle}
\le \norm{\mathbf{u}}\,\norm{\mathbf{v}},
\]
with equality exactly when the vectors are parallel. Now think of mean-zero random
variables as "vectors" and define their inner product to be
\[
\langle U, V \rangle \;=\; \E[UV].
\]
This genuinely \emph{is} an inner product: it is symmetric, linear in each slot, and
$\langle U, U\rangle = \E[U^2] \ge 0$, vanishing only when $U = 0$ (almost surely). Under this
dictionary,
\[
\langle U, V\rangle = \Cov(U, V), \qquad
\norm{U} = \sqrt{\E[U^2]} = \SD(U),
\]
so Cauchy--Schwarz reads
\[
\abs{\Cov(U, V)} \le \SD(U)\,\SD(V)
\quad\Longleftrightarrow\quad
\abs{\Corr(U, V)} \le 1.
\]
And the geometric meaning of the ratio is exactly the cosine of the angle $\theta$
between the two vectors:
\[
\Corr(X, Y) = \cos\theta = \frac{\langle U, V\rangle}{\norm{U}\,\norm{V}}.
\]
Uncorrelated variables ($\rho = 0$) are \emph{orthogonal} -- they meet at a right
angle. Perfectly correlated variables ($\rho = \pm 1$) are \emph{parallel} (or
anti-parallel). This is the same geometry as the angle between arrows in space; the
companion linear-algebra book develops the inner-product and Cauchy--Schwarz story
in full, and it provides the geometric foundation for correlation.
\end{connectionbox}

\begin{theorem}[Cauchy--Schwarz Inequality for Random Variables]
\label{ch14:thm-cauchy-schwarz}
For any random variables $X$ and $Y$ with finite second moments,
\[
\big(\Cov(X, Y)\big)^2 \;\le\; \Var(X)\,\Var(Y),
\]
with equality if and only if there exist constants $a, b$ (not both zero) such that $a(X - \mu_X) + b(Y - \mu_Y) = 0$ almost surely. Consequently, $-1 \le \Corr(X, Y) \le 1$.
\end{theorem}

\begin{proof}
Let $U = X - \mu_X$ and $V = Y - \mu_Y$ be the centered random variables. If $\Var(V) = 0$, then $V = 0$ almost surely, $\Cov(X,Y) = 0$, and the inequality holds immediately as $0 \le 0$. Thus assume $\Var(V) = \E[V^2] > 0$.

For any real parameter $t \in \R$, construct the random variable $U - tV$. Because a squared quantity is non-negative on every single outcome, $(U - tV)^2 \ge 0$ pointwise. Since expectation preserves non-negativity, the expected value must remain non-negative for every $t$:
\[
g(t) \;=\; \E\big[(U - tV)^2\big] \;\ge\; 0 \qquad \text{for all } t \in \R.
\]
Expanding the square using the linearity of expectation:
\begin{align*}
g(t) &= \E\big[U^2 - 2t\, UV + t^2 V^2\big] \\
     &= \E[V^2]\, t^2 \;-\; 2\E[UV]\, t \;+\; \E[U^2].
\end{align*}
This is a standard quadratic polynomial $A t^2 + B t + C$ in the variable $t$, whose coefficients are
\[
A = \E[V^2] = \Var(Y) > 0, \qquad B = -2\E[UV] = -2\Cov(X,Y), \qquad C = \E[U^2] = \Var(X).
\]
Because $g(t) \ge 0$ for all real $t$, this upward-opening parabola can never cross strictly below the horizontal axis. In algebra, a quadratic with positive leading coefficient $A > 0$ is non-negative everywhere if and only if its discriminant $\Delta = B^2 - 4AC$ is non-positive:
\[
\Delta \;=\; \big(-2\Cov(X,Y)\big)^2 \;-\; 4\,\Var(Y)\,\Var(X) \;\le\; 0.
\]
Simplifying the inequality:
\[
4\big(\Cov(X,Y)\big)^2 \;\le\; 4\,\Var(X)\,\Var(Y) \;\implies\; \big(\Cov(X,Y)\big)^2 \;\le\; \Var(X)\,\Var(Y).
\]
Taking square roots gives $|\Cov(X,Y)| \le \SD(X)\,\SD(Y)$, and dividing by $\SD(X)\,\SD(Y)$ establishes $-1 \le \Corr(X,Y) \le 1$.

Equality $\Delta = 0$ occurs if and only if the minimum of the parabola touches the axis at its vertex $t^\star = -B/(2A) = \Cov(X,Y)/\Var(Y)$. At that point, $\E[(U - t^\star V)^2] = 0$, which forces $U - t^\star V = 0$ almost surely, proving the exact linear relationship.
\end{proof}

Figure~\ref{ch14:fig-cosine} illustrates this geometry side by side. In the space of finite-variance random variables, centered random variables act as vectors, standard deviations give vector lengths, and the inner product is expectation.

\begin{figure}[H]
\centering
\begin{tikzpicture}[scale=0.92]
  % Panel 1: Hilbert space vector geometry
  \begin{scope}[shift={(0,0)}]
    \draw[draw=gray!30, fill=gray!3, rounded corners=4pt] (-0.8,-1.8) rectangle (5.6,4.0);
    \node[anchor=north, font=\small\bfseries, structurepurple!80!black] at (2.4,3.8) {1. Random Variables as Vectors ($L^2$ Space)};

    \coordinate (O) at (0,0);
    \coordinate (X) at (4.0,0);
    \coordinate (Y) at ({4.0*cos(42)}, {4.0*sin(42)});
    \coordinate (P) at ({4.0*cos(42)}, 0);

    % Base axis guide
    \draw[->, thick, gray!60] (-0.4,0) -- (4.8,0);

    % Vectors
    \draw[->, very thick, sourceblue] (O) -- (X);
    \node[anchor=south west, font=\footnotesize\bfseries, sourceblue] at ($(X)+(0.05, 0.05)$) {$\tilde X$ \scriptsize ($\|\tilde X\|=1$)};
    \draw[->, very thick, densityred] (O) -- (Y) node[above right=-2pt, font=\footnotesize\bfseries] {$\tilde Y$ \scriptsize ($\|\tilde Y\|=1$)};

    % Angle arc
    \draw[thick, gray!80] (1.0,0) arc (0:42:1.0);
    \node[font=\footnotesize\bfseries] at (21:1.3) {$\theta$};

    % Orthogonal projection
    \draw[dashed, structurepurple, thick] (Y) -- (P) node[midway, right, font=\tiny\bfseries, structurepurple] {$\tilde Y - \cos\theta\,\tilde X$};
    % Right angle symbol
    \draw[gray!70] ($(P)+(-0.2,0)$) -- ($(P)+(-0.2,0.2)$) -- ($(P)+(0,0.2)$);

    % Projection vector on X
    \draw[ultra thick, quotientorange] (O) -- (P);
    \fill[quotientorange] (P) circle (1.8pt);
    \node[below, font=\scriptsize\bfseries, quotientorange!90!black] at (P) {$\cos\theta\,\tilde X$};

    % Annotation
    \node[align=center, font=\scriptsize, draw=structurepurple!60, fill=white, rounded corners=2pt, inner sep=2pt]
      at (2.4,-1.1) {Inner Product: $\langle \tilde X, \tilde Y \rangle = \E[\tilde X \tilde Y] = \cos\theta = \Corr(X,Y)$\\[1pt]\scriptsize Orthogonal ($\theta = 90^\circ$) $\Longleftrightarrow$ Uncorrelated ($\Corr = 0$)};
  \end{scope}

  % Bridge arrow
  \draw[<->, >=stealth, very thick, gray!60] (5.8, 1.0) -- (7.0, 1.0)
    node[midway, above=2pt, font=\tiny\bfseries, align=center] {Geometry\\$\longleftrightarrow$\\Data};

  % Panel 2: Corresponding scatterplot in data space
  \begin{scope}[shift={(7.8,0)}]
    \draw[draw=gray!30, fill=gray!3, rounded corners=4pt] (-0.6,-1.8) rectangle (5.8,4.0);
    \node[anchor=north, font=\small\bfseries, targetgreen!80!black] at (2.6,3.8) {2. Scatter Cloud in Data Space};

    \draw[->, thick, gray!70] (0.2,1.0) -- (5.0,1.0) node[right, font=\scriptsize] {$\tilde X$};
    \draw[->, thick, gray!70] (2.6,-0.6) -- (2.6,2.8) node[above, font=\scriptsize] {$\tilde Y$};

    % Elliptical contour rotated by roughly 35 deg
    \draw[dashed, structurepurple!70, thick]
      (2.6,1.0) ellipse [x radius=1.9, y radius=0.75, rotate=35];

    % Regression line: slope = cos(42) ~ 0.74
    \draw[densityred, very thick] (0.8, -0.33) -- (4.4, 2.33)
      node[above right=-2pt, font=\tiny\bfseries, densityred] {$\hat Y = \rho\,\tilde X$};

    % Synthetic scatter points conforming to rho ~ 0.74
    \foreach \px/\py in {
      1.2/0.0, 1.5/0.4, 1.8/0.2, 2.0/0.8, 2.3/0.6,
      2.5/1.1, 2.7/0.9, 3.0/1.5, 3.2/1.3, 3.5/1.8,
      3.8/1.6, 4.0/2.2, 2.1/0.3, 3.1/1.1, 2.8/1.7
    }{
      \fill[sourceblue!80!black] (\px, \py) circle (1.4pt);
    }

    \node[align=center, font=\scriptsize, draw=targetgreen!80!black, fill=white, rounded corners=2pt, inner sep=2pt]
      at (2.6,-1.1) {Scatter Tilt: $\rho = \cos 42^\circ \approx 0.74$\\[1pt]\scriptsize Tighter ellipse $\implies$ smaller $\theta$ $\implies |\rho| \to 1$};
  \end{scope}
\end{tikzpicture}
\caption{Two geometric views of correlation. (Left) In the vector space of centered random variables (known formally as $L^2$ space), correlation is the cosine of the angle $\theta$ between unit vectors $\tilde X$ and $\tilde Y$, with orthogonal vectors corresponding to zero correlation. (Right) In data space, that same cosine $\rho = \cos\theta$ manifests as the slope of the standardized regression line and the eccentricity of the scatter ellipse.}
\label{ch14:fig-cosine}
\end{figure}

\begin{historybox}
The correlation coefficient was developed in the late nineteenth century through studies of biological inheritance. In the 1880s, \emph{Sir Francis Galton} set out to quantify how physical traits, like height and forearm length, pass from parents to children. He was stuck on a stubborn puzzle: how could one compare inheritance across different traits without the measurement units (inches versus centimeters) or differing population spreads distorting the comparison? 

Galton had few modern mathematical tools. He relied on hand-tallied cross-classification tables (counting how many parents of height $x$ had children of height $y$) and smoothed contour ellipses (ovals circling regions of equal data density on graph paper). He noticed that tall parents tended to have tall children who were, on average, slightly less extreme than themselves---a phenomenon he christened "regression toward mediocrity" (the birth of the word \emph{regression}). Galton recognized that a single scale-free number was hiding in the tilt of those ellipses, and he adapted the bivariate normal error equations derived in 1846 by French astronomer \emph{Auguste Bravais} (who had studied telescope sighting errors without realizing the equations measured co-variation). 

In 1895, Galton's prot\'eg\'e \emph{Karl Pearson} gave the concept mathematical precision, publishing the formula $\rho = \Cov / (\SD\cdot\SD)$ that we use today. Yet the breakthrough sparked fierce controversy. Galton's work was explicitly motivated by founding eugenics---a deeply troubling ideological doctrine that sought to control human reproduction and cast a long shadow over early British statistics. Conceptually, Pearson was an aggressive philosophical positivist (believing science should only report observed associations, never unseen causes); he dogmatically asserted that correlation replaced causation entirely in modern science. That dogmatism triggered bitter feuds with contemporaries like \emph{G. Udny Yule} and later \emph{Ronald A. Fisher}, who insisted that passive correlation can never substitute for causal mechanisms.

Pearson's correlation coefficient became a basic tool in linear regression, multivariate statistics, psychometrics, and econometrics. Today, the frontier is finding what Pearson's number misses: because $\rho$ captures only straight-line relationships, modern statistics actively develops tools like \emph{distance correlation} (which detects arbitrary nonlinear dependencies), copulas, and causal directed acyclic graphs to map relationships that curve or hide behind confounders.
\end{historybox}

\subsection{What different correlations look like}

Numbers like $\rho = 0.3$ or $\rho = 0.8$ only become meaningful once you have seen
the scatterplots they produce. Figure~\ref{ch14:fig-clouds} shows clouds at several
correlation levels: a correlation near $0$ forms a round cloud; as
$\abs{\rho}$ approaches $1$ the cloud contracts toward a line;
the sign of $\rho$ gives the direction of the tilt.

\begin{figure}[H]
\centering
\begin{tikzpicture}
\begin{axis}[
  width=5.6cm, height=5.6cm,
  title={$\rho \approx 0$},
  title style={font=\small},
  xmin=-3,xmax=3,ymin=-3,ymax=3, xtick=\empty, ytick=\empty,
  axis lines=box, enlargelimits=false,
]
\addplot[only marks, mark=*, mark size=0.7pt, draw=sourceblue, fill=sourceblue]
 table {
 0.4 -1.2
 -1.1 0.9
 1.6 0.2
 -0.5 -0.7
 0.9 1.4
 -1.8 -0.4
 0.2 1.9
 1.3 -1.5
 -0.9 1.6
 0.6 0.3
 -1.5 -1.7
 2.0 0.8
 -0.3 -1.9
 1.1 -0.6
 -2.1 1.1
 0.0 0.0
 0.7 -2.0
 -0.6 1.2
 1.8 -0.9
 -1.3 -0.2
 0.3 0.7
 -0.8 -1.1
 1.4 1.7
 -1.9 0.5
 0.5 -0.4
 };
\end{axis}
\end{tikzpicture}\hfill
\begin{tikzpicture}
\begin{axis}[
  width=5.6cm, height=5.6cm,
  title={$\rho \approx 0.5$},
  title style={font=\small},
  xmin=-3,xmax=3,ymin=-3,ymax=3, xtick=\empty, ytick=\empty,
  axis lines=box, enlargelimits=false,
]
\addplot[only marks, mark=*, mark size=0.7pt, draw=targetgreen, fill=targetgreen]
 table {
 -2.0 -1.3
 -1.5 -0.4
 -1.2 -1.6
 -0.8 0.2
 -0.6 -0.9
 -0.3 0.4
 -0.1 -0.6
 0.1 0.7
 0.3 -0.2
 0.5 1.0
 0.7 0.1
 0.9 1.3
 1.2 0.5
 1.5 1.6
 1.8 0.9
 2.1 1.9
 -1.7 0.3
 0.6 -1.0
 1.3 1.8
 -0.9 -1.4
 0.0 0.5
 1.0 0.0
 -0.4 -0.8
 1.6 0.7
 -1.1 -0.5
 };
\end{axis}
\end{tikzpicture}

\vspace{0.6em}
\begin{tikzpicture}
\begin{axis}[
  width=5.6cm, height=5.6cm,
  title={$\rho \approx 0.9$},
  title style={font=\small},
  xmin=-3,xmax=3,ymin=-3,ymax=3, xtick=\empty, ytick=\empty,
  axis lines=box, enlargelimits=false,
]
\addplot[only marks, mark=*, mark size=0.7pt, draw=quotientorange, fill=quotientorange]
 table {
 -2.1 -1.9
 -1.7 -1.5
 -1.4 -1.7
 -1.1 -0.8
 -0.8 -1.0
 -0.5 -0.3
 -0.3 -0.6
 -0.1 0.1
 0.1 -0.2
 0.4 0.6
 0.6 0.4
 0.9 1.1
 1.1 0.8
 1.4 1.6
 1.7 1.4
 2.0 2.1
 -1.9 -1.6
 0.7 0.9
 1.3 1.2
 -0.9 -1.2
 0.2 0.3
 1.0 0.7
 -0.6 -0.5
 1.6 1.8
 -1.2 -1.0
 };
\end{axis}
\end{tikzpicture}\hfill
\begin{tikzpicture}
\begin{axis}[
  width=5.6cm, height=5.6cm,
  title={$\rho \approx -0.8$},
  title style={font=\small},
  xmin=-3,xmax=3,ymin=-3,ymax=3, xtick=\empty, ytick=\empty,
  axis lines=box, enlargelimits=false,
]
\addplot[only marks, mark=*, mark size=0.7pt, draw=trapred, fill=trapred]
 table {
 -2.0 1.7
 -1.6 1.9
 -1.3 1.1
 -1.0 1.4
 -0.7 0.6
 -0.4 0.9
 -0.2 0.3
 0.0 0.0
 0.2 -0.4
 0.5 -0.7
 0.7 -0.3
 1.0 -1.1
 1.2 -0.8
 1.5 -1.6
 1.8 -1.4
 2.1 -1.9
 -1.8 1.3
 0.6 -0.9
 1.3 -1.2
 -0.9 1.0
 0.1 -0.2
 1.0 -0.6
 -0.5 0.7
 1.6 -1.7
 -1.2 1.2
 };
\end{axis}
\end{tikzpicture}
\caption{Scatter clouds at four correlation levels. Near $\rho = 0$ the cloud is
diffuse; as $\abs{\rho}\to 1$ it concentrates tightly along a line. Positive $\rho$
tilts up to the right; negative $\rho$ tilts down. Note that $\rho$ measures
\emph{tightness around a line}, not the line's steepness.}
\label{ch14:fig-clouds}
\end{figure}

\begin{trapbox}
\textbf{Correlation is not slope.} A common misreading is that a high correlation
means a steep relationship, but correlation actually measures how \emph{tightly} the
points cluster around \emph{some} line, regardless of how steep that line is. A
gently sloping but razor-thin cloud can have $\rho = 0.99$, while a steep but fat
cloud can have $\rho = 0.3$. Steepness is the job of the regression slope (a
different quantity); $\rho$ only reports the strength of the linear pattern, scaled
to $[-1,1]$.
\end{trapbox}

\subsection{Worked example: correlation of two dice}

\begin{example}[$\Corr(D_1, S)$]
\label{ch14:ex-corr-dice}
Return to Example~\ref{ch14:ex-dice}. We found $\Cov(D_1, S) = \tfrac{35}{12}$.
The variance of a single fair die is $\Var(D_1) = \tfrac{35}{12}$, and since the
two dice are independent, $\Var(S) = \Var(D_1) + \Var(D_2) = \tfrac{35}{6}$.
Therefore
\[
\Corr(D_1, S) = \frac{\Cov(D_1, S)}{\sqrt{\Var(D_1)\Var(S)}}
= \frac{35/12}{\sqrt{(35/12)(35/6)}}
= \frac{35/12}{(35/12)\sqrt{2}}
= \frac{1}{\sqrt{2}} \approx 0.707.
\]
A clean answer. The first die "explains" the sum, but only halfway -- the second
die is an equal, independent contributor, so $D_1$ accounts for one of two equal
sources of variation, and $\rho = 1/\sqrt 2$. (If you summed $n$ independent dice,
the correlation of any one die with the total would be $1/\sqrt n$, fading as the
crowd of other dice grows.)
\end{example}

\section{Correlation Versus Causation}
\label{ch14:causation}

Few distinctions in statistics matter more in everyday reasoning than the gap
between correlation and causation.

\begin{trapbox}
\textbf{Correlation is not causation.} A nonzero correlation between $X$ and $Y$
tells you they move together. It says \emph{nothing} about \emph{why}. The
co-movement could arise because $X$ causes $Y$, because $Y$ causes $X$, because some
third variable $Z$ drives both (a \emph{confounder}), or by sheer coincidence in a
finite sample. The formula cannot tell these apart, because it only ever saw the two
columns of numbers, never the machinery that generated them.
\end{trapbox}

\begin{paradoxbox}
\textbf{The ice cream and drowning paradox.} Across the year, ice cream sales and
drowning deaths are strongly positively correlated. Taken naively, this suggests ice
cream causes drowning -- ban the ice cream and save lives! Of course the real story
is a lurking third variable: \emph{summer}. Hot weather drives \emph{both} ice cream
sales \emph{and} swimming (hence drownings). Neither one causes the other; a
confounder, the season, drives both. This is the canonical illustration of a
\textbf{confounded} correlation, and the cure is to control for the confounder --
e.g. look at the correlation \emph{within} a fixed month, where it largely
disappears.

Search across enough unrelated time series and you will inevitably find pairs
that track each other by pure chance; the sample correlation is real as
arithmetic, even when nothing connects the two quantities.
\end{paradoxbox}

\begin{happeningbox}
Why can't a single number distinguish causation? Because causation is a statement
about \emph{interventions} -- what would happen to $Y$ if we reached in and
\emph{forced} $X$ to change -- whereas correlation is a statement about
\emph{observation} -- what $Y$ tends to do when we passively watch $X$ vary.
Observation and intervention coincide only when there are no confounders and no
reverse causation, which is precisely what we cannot verify from the data columns
alone. This is why the gold standard for establishing causation is the
\emph{randomized controlled experiment}: by assigning $X$ at random, the
experimenter severs any link between $X$ and lurking confounders, so that an
observed correlation can finally be read as causation (Figure~\ref{ch14:fig-causation}). The modern field of causal
inference (Rubin, Pearl) is built to recover causal claims from observational data
by modeling exactly the confounding structure that raw correlation ignores.
\end{happeningbox}

\begin{figure}[H]
\centering
\begin{tikzpicture}[
  >=Stealth,
  node distance=1.8cm,
  box/.style={draw=gray!50, fill=gray!5, rounded corners=4pt, font=\small\bfseries, align=center, inner sep=4pt},
  action/.style={draw=sourceblue, fill=sourceblue!10, rounded corners=4pt, font=\small\bfseries, align=center, inner sep=4pt},
  confound/.style={draw=trapred, fill=trapred!10, rounded corners=4pt, font=\small\bfseries, align=center, inner sep=4pt},
  outcome/.style={draw=targetgreen, fill=targetgreen!10, rounded corners=4pt, font=\small\bfseries, align=center, inner sep=4pt},
]
  % Panel 1: Confounded Observation
  \begin{scope}[shift={(0,0)}]
    \draw[draw=gray!30, fill=gray!3, rounded corners=6pt] (-3.1,-2.3) rectangle (3.1,1.9);
    \node[anchor=north, font=\footnotesize\bfseries, trapred!80!black] at (0,1.75) {1. Passive Observation (Confounded)};

    \node[confound] (Z1) at (0,0.5) {Confounder $Z$\\ \footnotesize(Summer Heat)};
    \node[box] (X1) at (-1.9,-1.2) {$X$\\ \footnotesize(Ice Cream)};
    \node[outcome] (Y1) at (1.9,-1.2) {$Y$\\ \footnotesize(Drownings)};

    \draw[->, very thick, trapred] (Z1) -- (X1);
    \draw[->, very thick, trapred] (Z1) -- (Y1);
    \draw[<->, dashed, very thick, red!70!black] (X1) -- (Y1)
      node[midway, above=2pt, font=\tiny\bfseries, align=center, text=red!70!black] {Spurious\\[-1pt]$\rho > 0$};
  \end{scope}

  % Panel 2: Randomized Intervention
  \begin{scope}[shift={(6.7,0)}]
    \draw[draw=gray!30, fill=gray!3, rounded corners=6pt] (-3.1,-2.3) rectangle (3.1,1.9);
    \node[anchor=north, font=\footnotesize\bfseries, targetgreen!80!black] at (0,1.75) {2. Randomized Trial (Intervention)};

    \node[confound] (Z2) at (0,0.5) {Confounder $Z$\\ \footnotesize(User Type)};
    \node[action] (X2) at (-1.9,-1.2) {Treatment $X$\\ \footnotesize(New Feature)};
    \node[outcome] (Y2) at (1.9,-1.2) {Outcome $Y$\\ \footnotesize(Retention)};

    % Randomizer node
    \node[draw=sourceblue, circle, fill=white, inner sep=2pt, font=\scriptsize\bfseries] (R) at (-2.0,0.4) {Flip};
    \draw[->, very thick, sourceblue] (R) -- (X2.north -| R) node[midway, left=2pt, font=\tiny\bfseries] {Random};

    % Severed arrow from Z to X
    \draw[->, dashed, thick, gray!60] (Z2) -- (X2);
    \node[font=\bfseries\large, trapred] at (-0.95,-0.4) {$\times$};
    \node[font=\tiny\bfseries, trapred, fill=gray!3, inner sep=1pt] at (-0.95,0.0) {severed};

    \draw[->, very thick, gray!60] (Z2) -- (Y2);
    \draw[->, very thick, targetgreen!80!black] (X2) -- (Y2)
      node[midway, above=2pt, font=\tiny\bfseries, align=center, text=targetgreen!80!black] {Causal\\[-1pt]Effect};
  \end{scope}
\end{tikzpicture}
\caption{Observation versus intervention. (Left) When a lurking confounder $Z$ drives both $X$ and $Y$, the variables move together and produce a positive correlation even with zero direct causal connection. (Right) Random assignment $R$ severs the link between confounder $Z$ and treatment $X$, ensuring that any remaining association between $X$ and $Y$ reflects a genuine causal effect.}
\label{ch14:fig-causation}
\end{figure}

\begin{applicationbox}
\textbf{A/B testing and why we randomize.} Suppose a company notices that users who
use its new feature have higher retention. Does the feature \emph{cause} retention?
Maybe -- or maybe the kind of person who tries new features was already more engaged
(a confounder: user type). The whole point of an \textbf{A/B test} is to assign the
feature at random. Randomization makes the treatment group and control group
statistically identical \emph{in expectation} on every confounder, observed or not,
so the only systematic difference is the feature itself. Then the observed
difference in retention -- a correlation between "got the feature" and "stayed" --
can be read causally. This is correlation-versus-causation turned into an
engineering discipline, and it underlies essentially all rigorous product
experimentation, clinical trials, and policy evaluation.
\end{applicationbox}

\section{The Covariance Matrix}
\label{ch14:cov-matrix}

So far we have correlated two variables at a time. Real data come with many columns
at once -- a person's height, weight, age, blood pressure; a portfolio's dozens of
asset returns. We want \emph{all} the pairwise covariances organized in one tidy
object. That object is the \textbf{covariance matrix}, and it is the bridge from
this chapter into the multivariate world.

\begin{definition}[Covariance matrix]
\label{ch14:def-cov-matrix}
For a random vector $\mathbf{X} = (X_1, X_2, \dots, X_n)^\top$, the
\textbf{covariance matrix} $\Sigma$ is the $n \times n$ matrix whose
$(i, j)$ entry is the covariance of the $i$-th and $j$-th components:
\[
\Sigma_{ij} = \Cov(X_i, X_j),
\qquad
\Sigma =
\begin{pmatrix}
\Var(X_1) & \Cov(X_1, X_2) & \cdots & \Cov(X_1, X_n) \\
\Cov(X_2, X_1) & \Var(X_2) & \cdots & \Cov(X_2, X_n) \\
\vdots & \vdots & \ddots & \vdots \\
\Cov(X_n, X_1) & \Cov(X_n, X_2) & \cdots & \Var(X_n)
\end{pmatrix}.
\]
In compact vector notation, with $\boldsymbol\mu = \E[\mathbf{X}]$,
\[
\Sigma = \E\big[(\mathbf{X} - \boldsymbol\mu)(\mathbf{X} - \boldsymbol\mu)^\top\big].
\]
\end{definition}

Look at the structure. The \emph{diagonal} entries are the individual variances --
each variable's covariance with itself. The \emph{off-diagonal} entries are the
pairwise covariances. The matrix is \textbf{symmetric} ($\Sigma_{ij} =
\Sigma_{ji}$) because covariance is symmetric. And the correlation matrix is the same
object with every entry $\Sigma_{ij}$ replaced by $\Sigma_{ij}/(\SD(X_i)\SD(X_j))$,
which puts $1$s down the diagonal.

\begin{proposition}[Covariance matrices are positive semidefinite]
\label{ch14:psd}
For any constant vector $\mathbf{a} \in \R^n$,
\[
\mathbf{a}^\top \Sigma\, \mathbf{a} = \Var\!\big(\mathbf{a}^\top \mathbf{X}\big)
= \Var\!\Big(\sum_i a_i X_i\Big) \;\ge\; 0.
\]
Hence $\Sigma$ is symmetric positive semidefinite.
\end{proposition}

\begin{whybox}
A matrix $\Sigma$ is called \textbf{positive semidefinite} if it never produces a negative value when sandwiched between any vector: $\mathbf{a}^\top \Sigma \mathbf{a} \ge 0$. The quadratic form $\mathbf{a}^\top \Sigma \mathbf{a}$ expands into
$\sum_{i,j} a_i a_j \Cov(X_i, X_j)$, which by bilinearity
(Theorem~\ref{ch14:var-sum}) is exactly $\Cov\big(\sum_i a_i X_i, \sum_j a_j
X_j\big) = \Var\big(\sum_i a_i X_i\big)$. A variance is never negative, so the form
is never negative. That is the entire proof.
\end{whybox}

\begin{connectionbox}
\textbf{From covariance matrix to PCA.} Because $\Sigma$ is symmetric and positive
semidefinite, the spectral theorem from linear algebra guarantees it has an
orthonormal basis of \emph{eigenvectors} with nonnegative \emph{eigenvalues}. That
fact is the engine of \textbf{Principal Component Analysis (PCA)}, the most widely
used dimensionality-reduction method in data science.

Here is the picture. Proposition~\ref{ch14:psd} says $\mathbf{a}^\top \Sigma
\mathbf{a}$ is the variance of the data projected onto direction $\mathbf{a}$. Ask:
\emph{which direction in space carries the most variance?} Maximizing
$\mathbf{a}^\top \Sigma \mathbf{a}$ over unit vectors $\mathbf{a}$ is a classic
eigenvalue problem, and its answer is the \emph{top eigenvector} of $\Sigma$ -- the
long axis of the data cloud. Its eigenvalue is the variance along that axis. The
next eigenvector gives the most-spread direction orthogonal to the first, and so on.
The eigenvectors are the natural axes of the ellipsoidal cloud, and the eigenvalues
say how stretched the cloud is along each. PCA simply rotates your data into these
axes and keeps the few with the largest eigenvalues, discarding directions where the
data barely varies. So the covariance of this chapter, stacked into a matrix and
handed to the spectral theorem, becomes the geometry of high-dimensional data. The
companion linear-algebra book develops eigenvectors, the spectral theorem, and
quadratic forms; here we see why they sit at the center of multivariate statistics.
\end{connectionbox}

\begin{figure}[H]
\centering
\begin{tikzpicture}
\begin{axis}[
  width=8.5cm, height=6.5cm,
  xlabel={$X_1$}, ylabel={$X_2$},
  xmin=-3.5,xmax=3.5,ymin=-3,ymax=3, xtick=\empty, ytick=\empty,
  axis lines=middle, enlargelimits=false,
]
\addplot[only marks, mark=*, mark size=0.8pt, draw=structurepurple, fill=structurepurple]
 table {
 -2.4 -1.5
 -1.9 -1.0
 -1.6 -1.3
 -1.2 -0.5
 -0.9 -0.9
 -0.6 -0.2
 -0.3 -0.5
 0.0 0.1
 0.2 -0.3
 0.5 0.5
 0.7 0.2
 1.0 0.8
 1.3 0.5
 1.6 1.1
 1.9 0.9
 2.3 1.4
 -2.0 -0.7
 0.6 -0.1
 1.4 1.3
 -1.1 -1.1
 0.1 0.4
 1.0 0.3
 -0.5 -0.8
 1.7 1.5
 -1.3 -0.6
 2.0 1.0
 -2.2 -1.2
 0.8 0.7
 };
% principal axis (top eigenvector direction)
\addplot[densityred, very thick, -{Stealth}, domain=0:2.7] ({x*0.86},{x*0.50});
\addplot[densityred, very thick, -{Stealth}, domain=0:2.7] ({-x*0.86},{-x*0.50});
\node[above right, densityred] at (axis cs:2.35,1.35) {\small PC1};
% second axis
\addplot[targetgreen, very thick, -{Stealth}, domain=0:1.1] ({-x*0.50},{x*0.86});
\node[targetgreen] at (axis cs:-0.9,1.5) {\small PC2};
\end{axis}
\end{tikzpicture}
\caption{A correlated cloud and its principal axes. The first principal component
(PC1, red) points along the direction of greatest variance -- the top eigenvector of
the covariance matrix $\Sigma$. The second (PC2, green) is orthogonal and captures
the leftover spread. PCA rewrites the data in these natural axes.}
\label{ch14:fig-pca}
\end{figure}

\begin{applicationbox}
\textbf{Portfolios and diversification.} In finance, let $R_1, \dots, R_n$ be the
returns of $n$ assets and let $\mathbf{w}$ be the fractions of your money in each
(so $\sum_i w_i = 1$). The portfolio return is $\mathbf{w}^\top \mathbf{R}$, and by
Proposition~\ref{ch14:psd} its risk (variance) is exactly
\[
\Var(\text{portfolio}) = \mathbf{w}^\top \Sigma\, \mathbf{w}.
\]
Diversification is now visible in the algebra. If the assets are weakly correlated,
the off-diagonal entries of $\Sigma$ are small, the cross terms barely add, and
spreading your money across many assets drives the portfolio variance down. If the
assets are strongly positively correlated, the off-diagonal terms are large, the
cross terms pile on, and "diversifying" buys you almost nothing -- everything sinks
together. Markowitz won a Nobel Prize for formalizing exactly this: choosing
$\mathbf{w}$ to minimize $\mathbf{w}^\top \Sigma \mathbf{w}$ for a desired expected
return. The entire theory of portfolio risk is the quadratic form of this chapter's
covariance matrix.
\end{applicationbox}

\section{Sample Covariance and Correlation}
\label{ch14:sample}

Everything above is about \emph{true} (population) covariance, defined with
expectations over the full distribution. In practice we never see the distribution;
we see a finite sample $(x_1, y_1), \dots, (x_n, y_n)$ and must \emph{estimate}.
The natural estimators replace expectations with averages.

\begin{definition}[Sample covariance and correlation]
\label{ch14:def-sample}
Given paired data with sample means $\bar x, \bar y$, the \textbf{sample covariance}
is
\[
s_{xy} = \frac{1}{n-1} \sum_{i=1}^n (x_i - \bar x)(y_i - \bar y),
\]
and the \textbf{sample correlation} (Pearson's $r$) is
\[
r = \frac{s_{xy}}{s_x\, s_y}
= \frac{\sum_{i=1}^n (x_i - \bar x)(y_i - \bar y)}
       {\sqrt{\sum_{i=1}^n (x_i - \bar x)^2}\,
        \sqrt{\sum_{i=1}^n (y_i - \bar y)^2}},
\]
where $s_x, s_y$ are the sample standard deviations.
\end{definition}

The divisor $n - 1$ rather than $n$ is the same \emph{Bessel correction} you met for
the sample variance. Because the sample means $\bar x$ and $\bar y$ are calculated from the sample itself, the observations sit slightly closer to $(\bar x, \bar y)$ than to the true, unknown population means $(\mu_X, \mu_Y)$. Measuring deviations from sample means systematically underestimates the true spread on average; dividing by $n - 1$ rather than $n$ exactly offsets this downward bias, making $s_{xy}$ an unbiased estimator of $\Cov(X, Y)$. (Conveniently, the $n - 1$ factor cancels completely in $r$ -- you get the exact same correlation whether you divide by $n$ or $n - 1$ top and bottom.)

\begin{example}[Sample correlation by hand]
\label{ch14:ex-sample}
Five paired observations:
\[
\begin{array}{c|ccccc}
x & 1 & 2 & 3 & 4 & 5 \\ \hline
y & 2 & 1 & 4 & 3 & 5
\end{array}
\]
Then $\bar x = 3$, $\bar y = 3$. Compute the deviation products:
\[
\begin{array}{c|c|c|c}
x_i - \bar x & y_i - \bar y & (x_i-\bar x)(y_i-\bar y) & (x_i-\bar x)^2 \\ \hline
-2 & -1 & 2 & 4 \\
-1 & -2 & 2 & 1 \\
 0 &  1 & 0 & 0 \\
 1 &  0 & 0 & 1 \\
 2 &  2 & 4 & 4
\end{array}
\]
The sums: $\sum (x_i-\bar x)(y_i-\bar y) = 8$, $\sum(x_i-\bar x)^2 = 10$, and by the
same arithmetic $\sum (y_i - \bar y)^2 = 1 + 4 + 1 + 0 + 4 = 10$. So
\[
r = \frac{8}{\sqrt{10}\,\sqrt{10}} = \frac{8}{10} = 0.8.
\]
A strong positive linear association -- the points rise together, with a little
scatter. If you plot them (Figure~\ref{ch14:fig-sample}), you will see the points
hugging the rising diagonal with one small swap (the pairs $(1,2),(2,1)$ and
$(3,4),(4,3)$ each cross), exactly the look of an $r = 0.8$ cloud.
\end{example}

\begin{figure}[H]
\centering
\begin{tikzpicture}
\begin{axis}[
  width=8cm, height=6cm,
  xlabel={$x$}, ylabel={$y$},
  xmin=0, xmax=6, ymin=0, ymax=6,
  xtick={1,2,3,4,5}, ytick={1,2,3,4,5},
  axis lines=middle, enlargelimits=false, clip=true,
]
% best-fit line: slope = s_xy / s_x^2 = 8/10 = 0.8, through (3,3): y = 0.8x + 0.6
\addplot[densityred, thick, domain=0.4:5.6] {0.8*x + 0.6};
\addplot[only marks, mark=*, mark size=2.6pt, draw=sourceblue, fill=sourceblue]
  coordinates {(1,2)(2,1)(3,4)(4,3)(5,5)};
\node[densityred] at (axis cs:4.6,1.5) {\small $y = 0.8x + 0.6$};
\end{axis}
\end{tikzpicture}
\caption{The five paired observations of Example~\ref{ch14:ex-sample} and their
least-squares line. The points cluster tightly around a rising line, which is what
$r = 0.8$ looks like.}
\label{ch14:fig-sample}
\end{figure}

\begin{labbox}
Take any dataset you can find with two numeric columns -- car weight vs. fuel
economy, study hours vs. grades, temperature vs. electricity use. (1) Make a
scatterplot first. Guess $r$ by eye before computing. (2) Compute $r$. Were you
close? (3) Now try a deliberately nonlinear pair, like $x$ and $x^2 - c$ over
a symmetric range, and watch $r$ collapse toward $0$ even though the relationship is
perfect. (4) Finally, find a pair with a sneaky confounder (e.g. shoe size and
reading ability across all ages of children) and watch a strong correlation
evaporate once you restrict to a single age. These four experiments encode the four
big lessons of this chapter.
\end{labbox}

\section{Chapter Summary}
\label{ch14:summary}

\begin{recoverybox}
\textbf{Chapter summary:}
\begin{itemize}
\item \textbf{Covariance} $\Cov(X,Y) = \E[(X-\mu_X)(Y-\mu_Y)] = \E[XY] - \E[X]\E[Y]$ measures the average product of deviations: it is positive when deviations tend to have the same sign, and negative when they have opposite signs.
\item It is \textbf{bilinear and symmetric}, with $\Cov(X,X) = \Var(X)$, which yields the master identity $\Var(X+Y) = \Var(X) + \Var(Y) + 2\Cov(X,Y)$.
\item \textbf{Independence} implies zero covariance, but the converse fails: covariance sees only straight-line co-movement and is blind to curved, symmetric dependencies like $Y = X^2$.
\item \textbf{Correlation} $\Corr(X,Y) = \Cov(X,Y)/(\SD(X)\SD(Y))$ standardizes covariance into a scale-free number in $[-1,1]$---the covariance of the $z$-scored variables.
\item The bound $\abs{\rho}\le 1$ is \textbf{Cauchy--Schwarz}: covariance is an inner product and correlation is the \emph{cosine of the angle} between centered variables. Hitting $\pm 1$ forces the scatter to collapse onto an exact line.
\item \textbf{Correlation is not causation}: confounders, reverse causation, and random noise produce spurious correlation. Controlled randomization breaks confounding to isolate causal mechanisms.
\item The \textbf{covariance matrix} $\Sigma$ gathers all pairwise covariances; it is symmetric positive semidefinite, and its eigenvectors form the principal axes of data variance (PCA).
\end{itemize}
\end{recoverybox}

Covariance and correlation summarize linear association with a single number.
When relationships are nonlinear, however, a single linear summary misses the structure.
The natural next step is to ask: given $X = x$, what is the expected value of $Y$?
Conditioning on $X$ leads to \textbf{conditional expectation} $\E[Y \given X]$,
which provides the best prediction of $Y$ given $X$ and forms the subject of Chapter~15.

\section{Exercises}
\label{ch14:exercises}

\begin{exercise}[Warm-up computation]
\label{ch14:ex-warmup}
Let $X$ and $Y$ have $\E[X] = 1$, $\E[Y] = 2$, $\E[XY] = 5$, $\Var(X) = 4$, and
$\Var(Y) = 9$. Find $\Cov(X,Y)$, $\Corr(X,Y)$, and $\Var(2X - Y + 3)$.
\end{exercise}

\begin{exercise}[Indicator covariance]
\label{ch14:ex-indicator}
Draw a single card from a standard 52-card deck. Let $A = \ind\{\text{card is a
heart}\}$ and $B = \ind\{\text{card is a queen}\}$. Compute $\Cov(A, B)$. Are $A$
and $B$ independent? Relate your two answers.
\end{exercise}

\begin{exercise}[Sign of covariance from a story]
\label{ch14:ex-sign}
For each pair, state whether you expect the covariance to be positive, negative, or
roughly zero, and explain in one sentence: (a) a city's daily high temperature and
its ice cream sales; (b) the age of a used car and its resale price; (c) a fair
coin's outcome ($0/1$) and an independent die's value; (d) a person's height and the
last digit of their phone number.
\end{exercise}

\begin{exercise}[Uncorrelated but dependent]
\label{ch14:ex-uncorr}
Let $\Theta$ be uniform on $[0, 2\pi)$ and set $X = \cos\Theta$, $Y = \sin\Theta$,
so $(X, Y)$ is a random point on the unit circle. Show that $\Cov(X, Y) = 0$ yet
$X$ and $Y$ are \emph{not} independent. (Hint for the second part: they satisfy
$X^2 + Y^2 = 1$ exactly.)
\end{exercise}

\begin{exercise}[Variance of a difference]
\label{ch14:ex-diff}
Two stocks have daily-return standard deviations $\SD(A) = 0.02$ and $\SD(B) =
0.03$, with correlation $\rho = 0.6$. You go "long $A$, short $B$" and hold the
portfolio $A - B$. Compute $\SD(A - B)$. Compare it to the case $\rho = -0.6$ and
explain, using Theorem~\ref{ch14:var-sum}, why negative correlation makes the spread
trade so much riskier here.
\end{exercise}

\begin{exercise}[Correlation is scale-invariant]
\label{ch14:ex-scale}
Prove from the definition that for constants $a, c > 0$ and any $b, d$,
$\Corr(aX + b, cY + d) = \Corr(X, Y)$. Then show that if $a < 0$ the correlation
flips sign. Where exactly in your derivation does the sign of $a$ enter?
\end{exercise}

\begin{exercise}[A correlation of exactly 1]
\label{ch14:ex-one}
Suppose $\Corr(X, Y) = 1$. Using the equality case of
Theorem~\ref{ch14:corr-bound} (or by computing $\Var(\tilde X - \tilde Y)$), show
that $Y = aX + b$ for some constants with $a > 0$, and express $a$ in terms of the
standard deviations. Interpret $a$ geometrically.
\end{exercise}

\begin{exercise}[Cauchy--Schwarz and the triangle inequality]
\label{ch14:ex-cs}
Theorem~\ref{ch14:thm-cauchy-schwarz} proved Cauchy--Schwarz by analyzing the quadratic $\E[(U - tV)^2] \ge 0$.
\begin{enumerate}[label=(\alph*)]
\item Use Cauchy--Schwarz to prove the \emph{triangle inequality for standard deviations}:
\[
\SD(X + Y) \;\le\; \SD(X) + \SD(Y).
\]
(Hint: expand $\Var(X+Y) = \Var(X) + \Var(Y) + 2\Cov(X,Y)$ and apply $\Cov(X,Y) \le \SD(X)\SD(Y)$.)
\item When does equality $\SD(X + Y) = \SD(X) + \SD(Y)$ hold? Interpret this condition geometrically in terms of the angle $\theta$ between the centered variables in Figure~\ref{ch14:fig-cosine}.
\end{enumerate}
\end{exercise}

\begin{exercise}[Confounding in numbers]
\label{ch14:ex-confound}
A confounder $Z$ is $+1$ or $-1$ with equal probability. Given $Z$, let
$X = Z + \eps_1$ and $Y = Z + \eps_2$ where $\eps_1, \eps_2$ are independent
$\Normal(0, 1)$ noise terms independent of $Z$. Compute $\Cov(X, Y)$ and
$\Corr(X, Y)$. Then argue that \emph{conditional on $Z$}, $X$ and $Y$ are
uncorrelated. Explain how this models the ice-cream-and-drowning paradox.
\end{exercise}

\begin{exercise}[A small covariance matrix]
\label{ch14:ex-matrix}
Let $\mathbf{X} = (X_1, X_2, X_3)$ have covariance matrix
\[
\Sigma = \begin{pmatrix} 4 & 2 & 0 \\ 2 & 9 & -3 \\ 0 & -3 & 4 \end{pmatrix}.
\]
(a) Read off $\Var(X_2)$ and $\Corr(X_2, X_3)$. (b) Compute $\Var(X_1 + X_2 + X_3)$
using $\mathbf{a}^\top \Sigma \mathbf{a}$ with $\mathbf{a} = (1,1,1)^\top$.
(c) Verify $\Sigma$ is symmetric, and check that the variance you found in (b) is
nonnegative, as positive semidefiniteness demands.
\end{exercise}

\begin{exercise}[Why divide by $n-1$? (conceptual)]
\label{ch14:ex-bessel}
Explain in words why the sample covariance uses the divisor $n - 1$ rather than $n$.
What goes wrong, in expectation, if you use $n$? Why does the choice of divisor not
affect the sample correlation $r$?
\end{exercise}
